\section{Bessel estimates with uniform remainders}\label{inf:sec-bessel}

The large-order estimates needed here are uniform as the ratio of
argument to order varies. We derive them from an outgoing integral
equation. The only asymptotic normalization used in that derivation is
the fixed-order, large-argument normalization of the Hankel function.
The recurrence identities and that normalization are as in
\cite[\S\S10.6, 10.17]{inf:dlmf}; the order derivative formula used
below is \cite[Eq.~10.21.17]{inf:dlmf}.

We write $Y_\nu$ for the Bessel function of the second kind and
$H^{(1)}_\nu=J_\nu+iY_\nu$ for the outgoing Hankel function.
The modified Bessel functions of the second kind are denoted by
$\mathrm K_\nu$.

\subsection{Zeros, radial norms, and differentiation in the order}

\begin{lemma}[Spacing and boundary normalization]\label{inf:bessel-spacing}
Let $\nu>1/2$. Every positive zero of $J_\nu$ is simple,
$j_{\nu,k}>\nu$, and
\begin{equation}\label{inf:bessel-spacing-bound}
 j_{\nu,k+1}-j_{\nu,k}>\pi.
\end{equation}
At a zero $z=j_{\nu,k}$ one has
\begin{equation}\label{inf:lommel}
 \int_0^1\zeta J_\nu(z\zeta)^2\,d\zeta
 =\frac12J_\nu'(z)^2.
\end{equation}
Consequently a radial Dirichlet eigenfunction, normalized in its radial
$L^2$ space and with its sign chosen appropriately, has boundary
normal derivative $\sqrt2\,z$.
\end{lemma}
\begin{proof}
A simultaneous zero of $J_\nu$ and $J'_\nu$ would give the zero solution
of its nonsingular ordinary differential equation at that point.
Let $v(x)=\sqrt{x}J_\nu(x)$. On an interval between two positive zeros,
\[
 \int |v'|^2=\int\left(1-\frac{\nu^2-1/4}{x^2}\right)|v|^2
 <\int|v|^2.
\]
The Dirichlet Poincar\'e inequality on an interval of length $L$ gives
$\pi^2/L^2<1$, proving \eqref{inf:bessel-spacing-bound}.
On $(0,1)$ the lowest eigenvalue of
$-\partial_\zeta^2+(\nu^2-1/4)/\zeta^2$ is at least
$\pi^2+\nu^2-1/4>\nu^2$, by the same Poincar\'e inequality and
$\zeta^{-2}\ge1$. Every eigenvalue is at least this lower bound, so
$j_{\nu,k}>\nu$.

For the integral identity put $w(\zeta,z)=J_\nu(z\zeta)$.
Subtract the equation for $w$ multiplied by $\partial_z w$ from its
$z$ derivative multiplied by $w$. Integration gives
\[
 [\zeta\{w(\partial_z w)_\zeta-w_\zeta\partial_z w\}]_0^1
 =-2z\int_0^1\zeta w^2\,d\zeta.
\]
At a zero $z$, $w(1,z)=0$, $w_\zeta(1,z)=zJ'_\nu(z)$ and
$\partial_z w(1,z)=J'_\nu(z)$. The term at zero vanishes by the
regular Bessel series. This proves \eqref{inf:lommel}. The radial
mass-unitary function is $\sqrt\zeta J_\nu(z\zeta)$, whose derivative
at one is $zJ'_\nu(z)$ and whose squared norm is the integral in
\eqref{inf:lommel}. Their ratio proves the last assertion.
\end{proof}

\begin{lemma}[Order derivative]\label{inf:bessel-order}
For $\nu\ge1/2$ and $x=j_{\nu,k}$,
\begin{equation}\label{inf:order-derivative}
 0\le \frac{\arccos(\nu/x)}{\sqrt{1-(\nu/x)^2}}
       -\frac{\partial j_{\nu,k}}{\partial\nu}
 \le\frac{\pi}{16x^2}.
\end{equation}
In particular,
\begin{equation}\label{inf:order-lipschitz}
 1-\frac\pi{16x^2}\le\frac{\partial j_{\nu,k}}{\partial\nu}
 \le\frac\pi2.
\end{equation}
The estimates are independent of $k$. If $x/\nu\to2$ and
$x\to\infty$, then
$\partial_\nu j_{\nu,k}\to2\pi/(3\sqrt3)$, also when $k$ varies.
\end{lemma}
\begin{proof}
Watson's order-derivative identity uses the modified Bessel function
$\mathrm K_0$ and reads
\[
 \frac{\partial j_{\nu,k}}{\partial\nu}
 =2x\int_0^\infty \mathrm K_0(2x\sinh v)e^{-2\nu v}\,dv.
\]
It applies to the simple positive zeros and positive orders in question.
At $\nu=1/2$ the regular Schr\"odinger solution is a multiple of
$\sin x$, so its positive zeros also satisfy $x>\nu$. For larger
orders this follows from Lemma~\ref{inf:bessel-spacing}.
Replacing $\sinh v$ by $v$ increases the integral. The replacement
integral is
\[
 2x\int_0^\infty \mathrm K_0(2xv)e^{-2\nu v}\,dv
 =\frac{\arccos(\nu/x)}{\sqrt{1-(\nu/x)^2}}.
\]
Insert the integral representation
\[
 \mathrm K_0(u)=\int_0^\infty e^{-u\cosh w}\,dw
\]
and integrate first in $v$. The result is
$\int_0^\infty(\cosh w+\nu/x)^{-1}\,dw$.
The substitution $z=\tanh(w/2)$ evaluates this integral to the
quotient above.

Since $-\mathrm K'_0=\mathrm K_1$ and $\mathrm K_1$ is positive and decreasing,
\[
 0\le \mathrm K_0(2xv)-\mathrm K_0(2x\sinh v)
 \le2x(\sinh v-v)\mathrm K_1(2xv).
\]
Taylor's integral formula gives $\sinh v-v\le v^3\cosh v/6$.
The hypothesis $\nu\ge1/2$ implies
$e^{-2\nu v}\cosh v\le1$. The error in the order derivative is
therefore at most
\[
 \frac{4x^2}{6}\int_0^\infty v^3\mathrm K_1(2xv)\,dv
 =\frac1{24x^2}\int_0^\infty u^3\mathrm K_1(u)\,du
 =\frac\pi{16x^2}.
\]
To evaluate the last integral, use
$\mathrm K_1(u)=\int_0^\infty e^{-u\cosh w}\cosh w\,dw$.
Tonelli's theorem and the elementary gamma integral give
$6\int_0^\infty\operatorname{sech}^3w\,dw=3\pi/2$.
The quotient in \eqref{inf:order-derivative}, written as
$\beta/\sin\beta$ for $0<\beta<\pi/2$, lies in $[1,\pi/2]$.
This proves \eqref{inf:order-lipschitz}. If $x/\nu\to2$ and
$x\to\infty$, the error $\pi/(16x^2)$ tends to zero, and the
quotient in \eqref{inf:order-derivative} tends to
$2\pi/(3\sqrt3)$, proving the limit.
\end{proof}

\subsection{The outgoing integral equation}

For $z>1$ define
\begin{align}
 \mathfrak s(z)&=\sqrt{z^2-1}-\arccos(1/z),
 &A_{\mathrm{Deb}}(z)&=(z^2-1)^{-1/4},\label{inf:debye-definitions}\\
 I(z)&=\frac1{4\sqrt{z^2-1}}+
                \frac5{12(z^2-1)^{3/2}},
 &V(z)&=\frac{z^2(z^2+4)}{4(z^2-1)^3}.\nonumber
\end{align}
The letters $I,V$ in this subsection denote scalar functions only.
The perturbation operator in the inverse argument will be denoted by
$\mathcal V$.

\begin{lemma}[First uniform correction]\label{inf:debye-first}
Let $\nu\ge1$, $z\ge3/2$, and set
\[
 \phi=\nu\mathfrak s(z)-\frac\pi4,\qquad
 \mathcal A_\nu(z)=\sqrt{\frac2{\pi\nu}}A_{\mathrm{Deb}}(z),
 \qquad b_\nu(z)=\frac{I(z)}{2\nu}.
\]
Then
\begin{align}
 \left|\frac{J_\nu(\nu z)}{\mathcal A_\nu(z)}
           -\cos\phi-b_\nu(z)\sin\phi\right|&\le\mathsf R_\nu^{\mathrm{Deb}}(z),
 \label{inf:debye-J}\\
 \left|\frac{Y_\nu(\nu z)}{\mathcal A_\nu(z)}
           -\sin\phi+b_\nu(z)\cos\phi\right|&\le\mathsf R_\nu^{\mathrm{Deb}}(z),
 \label{inf:debye-Y}\\
 \mathsf R_\nu^{\mathrm{Deb}}(z)&=
 \frac1{2\nu^2}\left\{V(z)+\frac{I(z)^2}{1-I(z)/\nu}\right\}
 <\frac6{5\nu^2}.\label{inf:debye-remainder}
\end{align}
Thus the same bound holds on every compact interval about $z=2$,
including a value of $z$ depending on $\nu$.
\end{lemma}
\begin{proof}
Put $\xi=\mathfrak s(z)$ and write
$H_\nu^{(1)}(\nu z)=A_{\mathrm{Deb}}(z)w(\xi)$.
Here $\mathfrak s'(z)=\sqrt{z^2-1}/z$ and
\[
 2\frac{A_{\mathrm{Deb}}'}{A_{\mathrm{Deb}}}
 +\frac{\mathfrak s''}{\mathfrak s'}+\frac1z=0,
 \qquad
 A_{\mathrm{Deb}}''+z^{-1}A_{\mathrm{Deb}}'
 =\frac{z^2+4}{4(z^2-1)^{9/4}}.
\]
Thus substitution in Bessel's equation cancels the first derivative
and gives
\[
 w_{\xi\xi}+\nu^2w=\psi(\xi)w,\qquad
 \psi(\mathfrak s(z))=-V(z).
\]
The outgoing Hankel normalization for fixed $\nu$ at infinity is
$w\sim\sqrt{2/(\pi\nu)}e^{i\nu\xi-i\pi/4}$, with the differentiated
normalization as well. This fixes the outgoing solution integrated
from infinity.
Define $g$ by
$w=\sqrt{2/(\pi\nu)}e^{i\nu\xi-i\pi/4}g$. Variation of constants yields
\begin{equation}\label{inf:volterra}
 g(\xi)=1+\frac1{2i\nu}\int_\xi^\infty
       \{e^{2i\nu(v-\xi)}-1\}\psi(v)g(v)\,dv.
\end{equation}
The potential is integrable. Indeed, with $u=\sqrt{z^2-1}$,
\begin{equation}\label{inf:potential-integral}
 \int_{\mathfrak s(z)}^\infty|\psi(v)|\,dv
 =\int_z^\infty\frac{y(y^2+4)}{4(y^2-1)^{5/2}}\,dy=I(z).
\end{equation}
The iterated integral of order $k$ in \eqref{inf:volterra} has norm at
most $(I(z)/\nu)^k/k!$. The series converges absolutely and gives the
unique bounded outgoing amplitude. Integrating first to a finite
endpoint and then using the Hankel normalization identifies that
solution with the given Hankel function. It follows that
$|g-1|\le e^{I/\nu}-1$.

The first iterate has nonoscillatory term $-iI/(2\nu)$. Its oscillatory
term has absolute value at most $V/(2\nu^2)$. To see this, integrate by
parts once and use
\[
 V'(z)=-\frac{z(z^4+10z^2+4)}{2(z^2-1)^4}<0,
 \qquad
 \int_\xi^\infty|\psi'(v)|\,dv=V(z).
\]
The contribution of the remaining iterates is bounded by
\[
 \frac1\nu\int_\xi^\infty |\psi(v)|
                       (e^{I(v)/\nu}-1)\,dv
 =e^{I/\nu}-1-I/\nu
 \le\frac{I^2}{2\nu^2(1-I/\nu)}.
\]
The final inequality follows by expanding the exponential and bounding
$1/k!\le1/2$ for $k\ge2$.
Thus $|g-(1-iI/(2\nu))|\le\mathsf R_\nu^{\mathrm{Deb}}$.
Taking real and imaginary parts after multiplication by $e^{i\phi}$
proves \eqref{inf:debye-J} and \eqref{inf:debye-Y}.

For $z\ge3/2$, $I(z)\le7/(6\sqrt5)<21/40$ and $V(z)\le9/5$.
In particular $I(z)/\nu<1$. Substitution gives
\[
 \mathsf R_\nu^{\mathrm{Deb}}(z)
 <\frac1{2\nu^2}\left(\frac95+
          \frac{(21/40)^2}{1-21/40}\right)
 =\frac{1809}{1520\nu^2}<\frac6{5\nu^2}.
\]
\end{proof}

\begin{lemma}[Any fixed order]\label{inf:debye-finite}
Let $z>1$, $\nu>0$, and $v=(z^2-1)^{-1/2}$. Define polynomials
$U_k(v)$ by $U_0=1$ and
\begin{equation}\label{inf:debye-recursion}
 U_{k+1}(v)=\frac{v^2(1+v^2)}2U_k'(v)
           +\frac18\int_0^v(1+5w^2)U_k(w)\,dw.
\end{equation}
For every fixed integer $L\ge0$, the amplitude in
\eqref{inf:volterra} satisfies
\begin{equation}\label{inf:debye-all-orders}
 \left|g-\sum_{k=0}^L\frac{U_k(v)}{(i\nu)^k}\right|
 \le\frac{2U_{L+1}(v)}{\nu^{L+1}}e^{I(z)/\nu}.
\end{equation}
If $I(z)/\nu<1$, the exponential can be replaced by
$(1-I(z)/\nu)^{-1}$. In particular,
\begin{equation}\label{inf:debye-polynomials}
 U_1(v)=\frac v8+\frac{5v^3}{24},\qquad
 U_2(v)=\frac{81v^2+462v^4+385v^6}{1152}.
\end{equation}
\end{lemma}
\begin{proof}
One has $dv/d\xi=-v^2(1+v^2)$ and
$V=v^2(1+v^2)(1+5v^2)/4$. The recursion is equivalent to
\[
 2\frac{dU_{k+1}}{d\xi}
 =\psi U_k-\frac{d^2U_k}{d\xi^2},\qquad U_{k+1}(\infty)=0.
\]
Every coefficient in \eqref{inf:debye-recursion} is nonnegative.
It follows that $U_k\ge0$, $dU_k/d\xi\le0$ and
$d^2U_k/d\xi^2\ge0$. Substitution of the truncated sum $g_L$ into
$g''+2i\nu g'-\psi g=0$ leaves only
$(U_L''-\psi U_L)/(i\nu)^L$. This numerator is nonnegative and
\[
 \int_\xi^\infty(U_L''+VU_L)\,d\xi=2U_{L+1}(v).
\]
The integral equation for $g-g_L$ has a forcing term bounded by
$2U_{L+1}/\nu^{L+1}$. In each iterated potential integral, its value
at a later point is no greater than its value at the initial point.
The same ordered-simplex estimate as for \eqref{inf:volterra}
therefore multiplies this bound by at most $e^{I/\nu}$. This proves
\eqref{inf:debye-all-orders}. The inequality
$e^w\le(1-w)^{-1}$ for $0\le w<1$ proves the alternative bound.
The two polynomials follow by integrating the recursion twice.
\end{proof}

\subsection{Fixed spectral windows}

\begin{lemma}[One radial state and no adjacent angular states]
\label{inf:window-isolation}
For each fixed $W>0$ there is $P_W$ such that for real $p\ge P_W$ and
$R\in[2p-1,2p+4]$, the set
\begin{equation}\label{inf:ball-window}
 \mathcal W_p(W)=\{(j,k):|z_{j,k}^2-R^2|\le W\}
\end{equation}
contains at most one state for each angular index $j$ and has no pair
of angular indices differing by one. The threshold is independent of
the split.
\end{lemma}
\begin{proof}
The argument interval has length $O_W(p^{-1})$, which is less than
$\pi$ for large $p$. Lemma~\ref{inf:bessel-spacing} proves the first
assertion. For the second we prove an explicit elementary separation.
Let $z\ge8$ be a zero of $J_\nu$, with $1/3\le\nu/z<1$, and set
$y(v)=J_{\nu+2}(z+v)/J_{\nu+1}(z)$. Consecutive orders have no common
positive zero, so this normalization is valid. The recurrences give
\[
 y(0)=\frac{2(\nu+1)}z\in[2/3,9/4],\qquad
 y'(0)=1-\frac{2(\nu+1)(\nu+2)}{z^2},\quad |y'(0)|\le2.
\]
For $|v|\le1/10$, both coefficients
$(z+v)^{-1}$ and $1-(\nu+2)^2/(z+v)^2$ of the differential equation
have absolute value below one. In fact
$(\nu+2)/(z-1/10)\le(z+2)/(z-1/10)\le100/79<\sqrt2$.
The first-order system for $(y,y')$ has maximum row norm at most two.
Its integral equation gives
$\max(|y(v)|,|y'(v)|)\le(9/4)e^{1/5}<3$. Therefore
$y(v)\ge2/3-3/10>0$ on this interval.

Two roots in \eqref{inf:ball-window} differ in argument by at most
$2W/R$ for large $p$, which is below $1/10$ after increasing $P_W$.
Their orders are at least $p-1$ and smaller than their roots.
Thus the smaller order and its root satisfy $1/3\le\nu/z<1$.
The separation just proved excludes adjacent angular indices.
\end{proof}

\begin{lemma}[Thin phase strip]\label{inf:phase-strip}
For each fixed $W>0$ there are $C_W,P_W<\infty$ with the following
property. Suppose $p\ge P_W$, $R\in[2p-1,2p+4]$, and
$p-1\le\nu\le R-p^{2/3}/2$. Define
\begin{align}
 \kappa(\nu,z)&=\sqrt{z^2-\nu^2},\nonumber\\
 \Xi(\nu,z)&=\kappa(\nu,z)-\nu\arccos(\nu/z)-\pi/4,\label{inf:phase-def}\\
 b(\nu,z)&=\frac1{8\kappa(\nu,z)}+
                       \frac{5\nu^2}{24\kappa(\nu,z)^3}.\nonumber
\end{align}
If $J_\nu(z)=0$ and $|z^2-R^2|\le W$, then
\begin{equation}\label{inf:phase-strip-bound}
 \dist\bigl(\Xi(\nu,R)-\arctan b(\nu,R),
                   \pi(\Z+1/2)\bigr)\le C_W/p.
\end{equation}
For $\nu=p-1+2j$, put
\begin{equation}\label{inf:phase-graph}
 \varphi_p(j)=\pi^{-1}\{\Xi(p-1+2j,R)
                    -\arctan b(p-1+2j,R)\}-\tfrac12.
\end{equation}
On a band $D_0\le R-\nu\le2D_0$ with
$D_0\ge p^{2/3}/2$ one has
\begin{equation}\label{inf:phase-curvature}
 \frac{c}{\sqrt{pD_0}}\le\varphi_p''(j)
 \le\frac{C}{\sqrt{pD_0}},
\end{equation}
where $c,C>0$ are absolute and the threshold may depend on $W$.
\end{lemma}
\begin{proof}
Write $D=R-\nu$. In the stated region and for
$|z-R|\le W/R$, one has $\nu\asymp p$, $z+\nu\asymp p$, and
$\kappa^2\asymp pD$, with absolute comparison constants. The case
$L=1$ of Lemma~\ref{inf:debye-finite} gives
\[
 \frac{J_\nu(z)}{\sqrt{2/\pi}\,\kappa^{-1/2}}
 =\cos\Xi+b\sin\Xi+e,\qquad
 |e|\le\frac{2U_2(\nu/\kappa)}{\nu^2(1-I(z/\nu)/\nu)}.
\]
Here
$I(z/\nu)/\nu=1/(4\kappa)+5\nu^2/(12\kappa^3)
=O(p^{1/2}D^{-3/2})=O(p^{-1/2})$.
The explicit polynomial $U_2$ gives
$|e|=O(pD^{-3})=O(p^{-1})$. Thus the denominator is bounded away
from zero uniformly. At a zero, divide by $\sqrt{1+b^2}$ and use
$\arcsin u\le2u$ for $0\le u\le1/2$ to obtain a distance
$O(p^{-1})$ from $\Xi-\arctan b$ to $\pi(\Z+1/2)$.
Now $\Xi_z=\kappa/z\le1$ and
\[
 b_z=-\frac z{8\kappa^3}-\frac{5\nu^2z}{8\kappa^5}=o(1).
\]
Replacing $z$ by $R$ changes the phase by $O_W(p^{-1})$, proving
\eqref{inf:phase-strip-bound}.

At fixed $R$, $\Xi_{\nu\nu}=1/\kappa$ and direct differentiation gives
\[
 b_\nu=\frac{13\nu}{24\kappa^3}+\frac{5\nu^3}{8\kappa^5},\qquad
 b_{\nu\nu}=\frac{13}{24\kappa^3}
               +\frac{7\nu^2}{2\kappa^5}
               +\frac{25\nu^4}{8\kappa^7}.
\]
It follows that
\[
 \kappa\,| (\arctan b)_{\nu\nu}|
 \le\kappa\{|b_{\nu\nu}|+2|b||b_\nu|^2\}
 =O(pD^{-3})+O(p^2D^{-6})=O(p^{-1}).
\]
The constants are absolute throughout the region. Therefore
$\varphi_p''=(4/\pi)\kappa^{-1}(1+O(p^{-1}))$.
Since $\kappa\asymp\sqrt{pD_0}$ on the band, this proves
\eqref{inf:phase-curvature}.
\end{proof}

\begin{lemma}[Convex lattice strip]\label{inf:lattice-strip}
Let $f\in C^2(I)$, where $I$ has length at most $L$, and suppose
$0<m_0\le f''\le M_0$. Let $\delta>0$ satisfy
\[
 \frac12M_0L^3+2\delta L<1.
\]
Then the integer points $(j,k)$ with $j\in I$ and
$|k-f(j)|\le\delta$ are collinear. If there is at most one such point
at each $j$, their number is at most
$2+4\sqrt{\delta/m_0}$, with an integer upper bound understood.
\end{lemma}
\begin{proof}
For three such points, the determinant with rows $(1,j,k)$ is an
integer. The determinant of the three graph values is
$(x_2-x_1)(x_3-x_1)(x_3-x_2)f[x_1,x_2,x_3]$ when the abscissae
are distinct. The integral formula for the second divided
difference bounds its absolute value by $M_0/2$, so this
contribution is at most $M_0L^3/2$. If two abscissae coincide,
the graph determinant is zero and the same bound holds. The three vertical errors contribute at most
$2\delta L$: for ordered abscissae their coefficient absolute values
sum to twice the largest separation. Thus the determinant has
absolute value less than one and is zero. If there are at least three
points, fixing the two extreme ones shows that every point lies on
their line.

Take a middle point with abscissa $x_2$ and extreme abscissae
$x_1,x_3$. The chord of $f$ satisfies
\[
 \operatorname{chord}(x_2)-f(x_2)
 \ge\frac{m_0}{2}(x_2-x_1)(x_3-x_2).
\]
One proof subtracts $m_0x^2/2$ from $f$ and uses convexity of the
result. Since the three integer points are collinear, their errors
bound the same chord deficit by $2\delta$. Among $n$ distinct integer
abscissae choose a middle one; its two separations have product at
least $(n-2)^2/4$. Therefore $(n-2)^2\le16\delta/m_0$.
The assertions for zero, one or two points need no determinant.
\end{proof}

\begin{theorem}[Sublinear window count]\label{inf:global-count}
For every fixed $W>0$ there are $C_W,P_W<\infty$ such that, for
$p\ge P_W$ and $R\in[2p-1,2p+4]$,
\begin{equation}\label{inf:global-count-bound}
 \#\mathcal W_p(W)\le C_Wp^{2/3}.
\end{equation}
The bound holds for real $p$, and is independent of the split.
\end{theorem}
\begin{proof}
There is at most one radial state per angular index by
Lemma~\ref{inf:window-isolation}. Also $z_{j,k}>\nu_j$, so the
region $\nu_j>R-p^{2/3}$ has only $O_W(p^{2/3})$ possible indices.
In the remaining region Lemma~\ref{inf:phase-strip} assigns to each
state a lattice point $(j,m)$ with
$|m-\varphi_p(j)|\le\delta=C_W'/p$. This assignment is injective:
there is at most one radial state at a given $j$, and $2\delta<1$.

Divide $D=R-\nu_j$ into dyadic bands $D_0\le D\le2D_0$, from
$p^{2/3}$ up to $O(p)$. On each band
$f''\asymp(pD_0)^{-1/2}$ for $f=\varphi_p$. Divide its horizontal
interval into intervals of length
$L=c_*(pD_0)^{1/6}$, with $c_*>0$ fixed so small that
$M_0L^3/2\le1/4$. Since $D_0=O(p)$, the other term
$2\delta L=O_W(p^{-2/3})$ is below $1/4$ for large $p$.
Lemma~\ref{inf:lattice-strip} bounds the number of points per interval
by an absolute multiple of $1+\sqrt{C'_W}$, because
$\delta/m_0=O_W(\sqrt{D_0/p})=O_W(1)$.

The band has horizontal length $O(D_0)$, so it requires
$O(p^{-1/6}D_0^{5/6}+1)$ such intervals. The sum of $D_0^{5/6}$ over
the dyadic bands is $O(p^{5/6})$; the sum of the additive ones is
$O(\log p)$. This proves \eqref{inf:global-count-bound} after adding
the turning region. Every constant and threshold in this proof is
fixed once $W$ is fixed.
\end{proof}

\begin{corollary}[Domains in a thin shell]\label{inf:shell-count}
Fix $W,C_{\mathrm{shell}}>0$ and integer blocks $a,b$ with
$a+b=2p$ in the admissible range. If a two-block invariant domain in the unit-ball
normalization satisfies
$B_{1-C_{\mathrm{shell}}/p^2}\subset\Omega\subset B_{1+C_{\mathrm{shell}}/p^2}$, then its invariant
Dirichlet spectrum has at most $C_{W,C_{\mathrm{shell}}}p^{2/3}$ eigenvalues in
$[R^2-W,R^2+W]$ for large $p$ and $R\in[2p-1,2p+4]$.
The constants are uniform in the admissible integer splits.
\end{corollary}
\begin{proof}
Zero extension preserves invariance. The min--max principle and radial
dilation give, for each ordered eigenvalue index $i$,
\[
 (1+C_{\mathrm{shell}}/p^2)^{-2}\lambda_i(B)
 \le\lambda_i(\Omega)\le
 (1-C_{\mathrm{shell}}/p^2)^{-2}\lambda_i(B).
\]
Any index counted for $\Omega$ has its ball eigenvalue in
$[R^2-W',R^2+W']$, where $W'$ depends only on $W,C_{\mathrm{shell}}$ since
$R^2=O(p^2)$. Apply Theorem~\ref{inf:global-count}.
\end{proof}
